\documentclass[10.5pt,a4paper]{article}
\usepackage{geometry}
\geometry{margin=18mm,top=20mm,bottom=20mm}
\usepackage{kotex}
\usepackage{fontspec}
\setmainhangulfont{Apple SD Gothic Neo}[BoldFont={* Bold}]
\setsanshangulfont{Apple SD Gothic Neo}[BoldFont={* Bold}]
\setmonofont{Menlo}[Scale=0.82]
\usepackage{amsmath,amssymb}
\usepackage{booktabs,longtable,array}
\usepackage{graphicx}
\usepackage{enumitem}
\usepackage{xcolor}
\usepackage{fancyhdr}
\usepackage[hidelinks]{hyperref}
\hypersetup{colorlinks=true,linkcolor=blue!60!black,urlcolor=blue!65!black}
\setlength{\parindent}{0pt}
\setlength{\parskip}{4pt}
\renewcommand{\arraystretch}{1.18}
\setcounter{secnumdepth}{0}
\sloppy
\pagestyle{fancy}
\fancyhf{}
\fancyhead[L]{\small\color{gray}시스템 아키텍처: 카메라 중심 시계열 자차 중심 World Model}
\fancyhead[R]{\small\color{gray}KITTI Tracking Temporal World Model}
\fancyfoot[C]{\small\thepage}
\begin{document}
\begin{center}
{\LARGE\bfseries 시스템 아키텍처: 카메라 중심 시계열 자차 중심 World Model}\\[6pt]
{\normalsize\color{gray} 시스템 아키텍처: 카메라 중심 시계열 자차 중심 World Model}\\[4pt]
{\small\href{https://github.com/Heisnotanimposter/Totalcare_autonomous_driving_Yolov5}{https://github.com/Heisnotanimposter/Totalcare\_autonomous\_driving\_Yolov5}}
\end{center}
\vspace{4pt}
\hrule
\vspace{8pt}

\section{1. 시스템 개요}

본 프로젝트는 단일 전방 카메라 영상을 입력받아 구조화된 \textbf{시계열 자차 중심 World Model(Temporal Ego-centric World Model)}, 다중 시점 미래 궤적 예측, 위험도 지표, 그리고 권고형 Safety Shield 출력을 생성하는 7단계 파이프라인을 구현합니다.


\begin{quote}\begin{footnotesize}\begin{verbatim}
[Stage 1: 입력 영상 (KITTI 10 Hz)]
        |
        v
[Stage 2: 시각 인식 (src/kitti_prep.py, YOLO11m/YOLO11l + MobileSAM + SegFormer-B2)]
        |  2D BBox, 클래스(Car, Pedestrian, Cyclist), 검출 신뢰도
        |  희소 클래스 오버샘플링 + 인스턴스 폴리곤 + 도로/비도로/자차로 세그멘테이션
        v
[Stage 3: 시계열 이해 (src/track_utils.py, ByteTrack + MobileSAM)]
        |  시퀀스 격리 Track ID, 프레임 ID, 타임스탬프, mask_area_px, mask_fill_ratio, mask_bottom_xy
        v
[Stage 4: 자차 중심 World Model (src/world_model.py, src/road_lane.py)]
        |  P2 역분산 깊이 융합 (평면 지면 + 클래스 높이 사전값 + 마스크 최하단 접지점)
        |  OXTS IMU/GPS 자차 움직임 보상 (EgoPose2D: Ego <-> World ENU)
        |  상태 머신: Past <-> Present <-> Future | Observed / Inferred / Occluded / Unknown
        v
[Stage 5: 미래 궤적 예측 (src/world_model.py, World-ENU Kalman + CTRV Ego)]
        |  t + 0.5s, 1.0s, 1.5s, 2.0s 평균 위치 + 신뢰도/마스크 기반 1-sigma 공분산 타원
        v
[Stage 6: 위험도 평가 (src/world_model.py, C1 포텐셜 + CollisionRiskSmoother)]
        |  부드러운 시그모이드 위험 포텐셜 + 비대칭 EMA + 슬루율 |dR/dt| 제한 + 보정 TTC
        v
[Stage 7: 안전 인지 제어 권고 (src/world_model.py & src/world_visualizer.py)]
           히스테리시스 디바운싱 권고형 Safety Shield + 좌우 분할 화면(마스크 + BEV World)
\end{verbatim}\end{footnotesize}\end{quote}

\par\noindent\rule{\textwidth}{0.4pt}\par\medskip

\section{2. 좌표계 및 변환 수식}

\subsection{2.1 KITTI 카메라 좌표계}
KITTI Tracking (\texttt{calib/\allowbreak{}XXXX.\allowbreak{}txt}, \texttt{label\_\allowbreak{}02/\allowbreak{}XXXX.\allowbreak{}txt}) 기준:

\begin{itemize}[leftmargin=1.5em,itemsep=2pt]
\item $X_{\text{cam}}$: 우측 방향 ($+$)
\item $Y_{\text{cam}}$: 아래 방향 ($+$)
\item $Z_{\text{cam}}$: 광학축 전방 ($+$)
\end{itemize}

\subsection{2.2 자차 중심 BEV 좌표계}
\texttt{KittiCalibration} (\href{https://github.com/Heisnotanimposter/Totalcare\_autonomous\_driving\_Yolov5/blob/main/src/world\_model.py}{src/world\_model.py}) 에 정의됨:

\begin{itemize}[leftmargin=1.5em,itemsep=2pt]
\item $X_{\text{ego}}$: 횡방향 우측 ($+$), 단위 m
\item $Y_{\text{ego}}$: 종방향 전방 ($+$), 단위 m
\item $Z_{\text{ego}}$: 수직 상방 ($+$), 단위 m
\end{itemize}

표준 수평 장착($h_{\text{cam}} = 1.65\text{ m}$, 피치/롤 0 가정) 기준:

\[X_{\text{ego}} = X_{\text{cam}}, \quad Y_{\text{ego}} = Z_{\text{cam}}, \quad Z_{\text{ego}} = h_{\text{cam}} - Y_{\text{cam}}\]

\subsection{2.3 보정된 역분산 깊이 융합 및 해석적 불확실성}
KITTI 내부 파라미터 $P_2$ ($f_x, f_y, c_x, c_y$), 지면 접지점 $(u_{\text{mid}}, v_{\text{bottom}})$ (\texttt{mask\_\allowbreak{}fill\_\allowbreak{}ratio >= 0.\allowbreak{}25}일 때 MobileSAM 전경 최하단 \texttt{mask\_\allowbreak{}bottom\_\allowbreak{}xy}, 그 외 2D 박스 하단 중심 사용), 그리고 클래스 표준 높이 사전값 $H_c$ (\texttt{CLASS\_\allowbreak{}HEIGHT\_\allowbreak{}PRIORS\_\allowbreak{}M}: \texttt{Car=1.\allowbreak{}52m}, \texttt{Pedestrian=1.\allowbreak{}72m}, \texttt{Cyclist=1.\allowbreak{}68m})가 주어질 때:

\begin{enumerate}[leftmargin=1.8em,itemsep=2pt]
\item \textbf{평면 지면 투영 ($Y_{\text{ground}}$)} 및 피치 불확실성 $\sigma_\theta = 0.012\text{ rad}$:
\end{enumerate}
\[Y_{\text{ground}} = \frac{f_y \cdot h_{\text{cam}}}{v_{\text{bottom}} - c_y}, \qquad \sigma_{\text{ground}} = \sqrt{\left(\frac{Y_{\text{ground}}^2}{f_y h_{\text{cam}}}\sigma_v\right)^2 + \left(\frac{Y_{\text{ground}}^2}{h_{\text{cam}}}\sigma_\theta\right)^2}\]
\begin{enumerate}[leftmargin=1.8em,itemsep=2pt]
\item \textbf{피치 불변 높이 사전 깊이 ($Y_{\text{height}}$)} (화면 경계 미절단 박스 대상):
\end{enumerate}
\[Y_{\text{height}} = \frac{f_y \cdot H_c}{h_{\text{px}}}, \qquad X_{\text{ego}} = \frac{(u_{\text{mid}} - c_x) \cdot Y_{\text{fused}}}{f_x}\]
\begin{enumerate}[leftmargin=1.8em,itemsep=2pt]
\item \textbf{역분산 융합 (Inverse-Variance Fusion)}:
\end{enumerate}
\[w_g = \sigma_{\text{ground}}^{-2}, \quad w_h = \sigma_{\text{height}}^{-2}, \qquad Y_{\text{fused}} = \frac{w_g Y_{\text{ground}} + w_h Y_{\text{height}}}{w_g + w_h}, \qquad \sigma_Y = (w_g + w_h)^{-1/2}\]
World-ENU 칼만 필터링의 측정 공분산 $R_w$는 검출 신뢰도와 MobileSAM 마스크 채움 품질(\texttt{conf}, \texttt{mask\_\allowbreak{}fill\_\allowbreak{}ratio})에 반비례하여 추가 스케일링됩니다.


\textbf{무결성 원칙:}

\begin{itemize}[leftmargin=1.5em,itemsep=2pt]
\item $P_2$ 캘리브레이션이 누락된 경우 \texttt{TemporalWorldModel}은 \texttt{calibrated = False}로 설정하고 미터법 속도(\texttt{m/\allowbreak{}s}) 및 충돌 여유 시간(\texttt{TTC}) 산출을 \textbf{거부}합니다.
\item 투영된 3차원 위치는 항상 \textbf{추정값(Inferred)}(단안 지면/높이 융합 추정치)으로 분류되며 \textbf{관측값(Observed)} 으로 표기하지 않습니다.
\end{itemize}

\par\noindent\rule{\textwidth}{0.4pt}\par\medskip

\section{3. 자차 움직임 보상 및 시계열 World Model (\texttt{과거 $\rightarrow$ 현재 $\rightarrow$ 미래})}

\texttt{EgoPose2D}, \texttt{KittiOxtsOdometry}, \texttt{TemporalWorldModel}, \texttt{ObjectWorldTrack} (\href{https://github.com/Heisnotanimposter/Totalcare\_autonomous\_driving\_Yolov5/blob/main/src/world\_model.py}{src/world\_model.py})에 구현됨:


\begin{enumerate}[leftmargin=1.8em,itemsep=2pt]
\item \textbf{OXTS 자차 움직임 보상 (\texttt{Ego $\leftrightarrow$ World ENU})}:
\begin{itemize}[leftmargin=1.5em,itemsep=2pt]
\item \texttt{KittiOxtsOdometry}는 KITTI \texttt{oxts/\allowbreak{}<seq>.\allowbreak{}txt}의 \texttt{(lat,\allowbreak{} lon,\allowbreak{} yaw,\allowbreak{} vf,\allowbreak{} vl,\allowbreak{} wz)}를 0번 프레임 기준 메르카토르 투영을 통해 지역 미터법 \textbf{동-북(ENU)} 좌표계로 변환합니다.
\item \texttt{EgoPose2D.\allowbreak{}ego\_\allowbreak{}to\_\allowbreak{}world\_\allowbreak{}xy}와 \texttt{world\_\allowbreak{}to\_\allowbreak{}ego\_\allowbreak{}xy}는 순간 자차 중심 검출 좌표(\texttt{+X} 우측, \texttt{+Y} 전방)를 정지된 World-ENU 좌표계로 상호 변환합니다(실측 왕복 오차 \texttt{< 4e-\allowbreak{}14 m}).
\item 객체 속도와 4D 칼만 상태 \texttt{[X\_\allowbreak{}w,\allowbreak{} Y\_\allowbreak{}w,\allowbreak{} vX\_\allowbreak{}w,\allowbreak{} vY\_\allowbreak{}w]}는 \textbf{World-ENU}에서 필터링(자차 병진·요 회전 제거)되며, 자차의 계획 주행 경로는 \texttt{vf}와 \texttt{wz}를 이용한 등선회율·등속(\texttt{predict\_\allowbreak{}ctrv}) 롤아웃으로 산출됩니다.
\end{itemize}
\end{enumerate}

\begin{enumerate}[leftmargin=1.8em,itemsep=2pt]
\item \textbf{시계열 상태 전이 (\texttt{이전 미래 $\rightarrow$ 현재 $\rightarrow$ 과거})}:
\begin{itemize}[leftmargin=1.5em,itemsep=2pt]
\item 시각 $t_k$에서 활성 트랙의 World 좌표계 4D 운동학 칼만 상태와 공분산 $P$를 $\Delta t = t_k - t_{k-1}$만큼 전파한 뒤, 시각 $t_k$의 자차 BEV 좌표계로 재투영합니다.
\item 현재 프레임 검출과 매칭된 트랙은 회전된 측정 공분산 $R_w = R_{\text{ego}\to w} \, \text{diag}(\sigma_X^2, \sigma_Y^2) \, R_{\text{ego}\to w}^T$로 갱신되고 \texttt{observation\_\allowbreak{}state = "Observed"}를 부여받습니다.
\item 시각 $t_k$에 미관측되었으나 \texttt{max\_\allowbreak{}coast\_\allowbreak{}sec} ($1.0\text{ s}$) 이내인 트랙은 지수 신뢰도 감쇠 $\exp(-\Delta t_{\text{miss}} / 0.6)$와 함께 \texttt{observation\_\allowbreak{}state = "Occluded"}로 전이됩니다.
\item \texttt{max\_\allowbreak{}coast\_\allowbreak{}sec}를 초과하여 미관측된 트랙은 \texttt{last\_\allowbreak{}seen\_\allowbreak{}ts} 및 World 기준 위치 이력과 함께 \texttt{past\_\allowbreak{}world.\allowbreak{}expired\_\allowbreak{}tracks}로 보관됩니다.
\item \texttt{과거(Past)}, \texttt{현재(Present)}, \texttt{미래(Future)}는 시각 $t_k$ 기준 상대 타임스탬프로 색인되고 World-ENU에서 유지되므로, 자차 회전이나 후진 상황에서도 시계열 의미가 불변입니다.
\end{itemize}
\end{enumerate}

\begin{enumerate}[leftmargin=1.8em,itemsep=2pt]
\item \textbf{공간 정보 상태 (\texttt{Observed /\allowbreak{} Inferred /\allowbreak{} Occluded /\allowbreak{} Unknown})}:
\begin{itemize}[leftmargin=1.5em,itemsep=2pt]
\item \textbf{Observed (관측)}: 전방 카메라 화각(FOV) 내에서 현재 프레임에 직접 검출·추적된 대상.
\item \textbf{Inferred (추정)}: 지면 투영 BEV 위치, 운동학 속도, 칼만 미래 궤적(\texttt{+0.\allowbreak{}5s}\textasciitilde{}\texttt{+2.\allowbreak{}0s}), 위험도 지표.
\item \textbf{Occluded (폐색)}: 전경 객체 후방으로 레이캐스트된 각도 음영 영역(\texttt{compute\_\allowbreak{}occlusion\_\allowbreak{}sectors}) 및 일시 유실 유지 트랙.
\item \textbf{Unknown (미관측 영역)}: 자차 후방($Y_{\text{ego}} < 0$) 및 전방 카메라 수평 화각(\texttt{2 * atan2(W /\allowbreak{} 2,\allowbreak{} fx)}) 외부 영역.
\end{itemize}
\end{enumerate}

\par\noindent\rule{\textwidth}{0.4pt}\par\medskip

\section{4. Safety Shield 정책 및 충돌 감지기 기울기 평활화 (권고 전용)}

\texttt{assess\_\allowbreak{}risk\_\allowbreak{}and\_\allowbreak{}safety\_\allowbreak{}shield}는 다음 요소를 결합합니다:

\begin{enumerate}[leftmargin=1.8em,itemsep=2pt]
\item \textbf{$C^1$ 연속 미분 가능 위험 포텐셜}: 개별 객체 위험 점수는 실제 \textbf{횡방향 공분산}(\texttt{lateral\_\allowbreak{}uncertainty\_\allowbreak{}m}, \texttt{cov\_\allowbreak{}future\_\allowbreak{}ego[0][0]}) 기반 차로 침범 판정, 접근 속도 게이팅(\texttt{approach\_\allowbreak{}w}, 정지 차량·보도 오경보 방지), 보정 TTC, 그리고 소형 최조 충돌 구간 \texttt{conflict\_\allowbreak{}zones} (\texttt{[1.\allowbreak{}0,\allowbreak{} 2.\allowbreak{}2] m}, 트랙당 최대 1개)에 부드러운 \texttt{\_\allowbreak{}sigmoid} 곡선을 적용하며, 프레임 전체 위험도는 확률적 soft-union $1 - \prod_i (1 - r_i)$로 집계됩니다.
\item \textbf{종횡비 보정 BEV 불확실성 타원 (\texttt{BevCanvasMapper.\allowbreak{}cov\_\allowbreak{}to\_\allowbreak{}pixel\_\allowbreak{}ellipse})}: World-ENU $2\times 2$ 공분산을 야코비안 $J = \text{diag}(s_x, -s_y)$로 BEV 픽셀 좌표에 투영하여 종방향 깊이 불확실성이 인접 차로로 과대 확장되지 않도록 하며, 반투명 시점 채움(\texttt{t = +1.\allowbreak{}0s,\allowbreak{} +2.\allowbreak{}0s})과 단일 \texttt{+2.\allowbreak{}0s} 외곽선으로 시각화합니다.
\item \textbf{기울기 평활화 (\texttt{CollisionRiskSmoother})}: 비대칭 EMA(\texttt{alpha\_\allowbreak{}up=0.\allowbreak{}36}, \texttt{alpha\_\allowbreak{}down=0.\allowbreak{}22}), 명시적 슬루율 제한(\texttt{max\_\allowbreak{}rise\_\allowbreak{}per\_\allowbreak{}sec=1.\allowbreak{}10/\allowbreak{}s}, \texttt{max\_\allowbreak{}fall\_\allowbreak{}per\_\allowbreak{}sec=0.\allowbreak{}85/\allowbreak{}s}), 트랙별 위험도 EMA(\texttt{per\_\allowbreak{}track\_\allowbreak{}ema}), 저-TTC 연속 프레임 디바운싱, 그리고 Safety Shield 임계값 경계 히스테리시스(\texttt{hysteresis\_\allowbreak{}margin=0.\allowbreak{}04})를 통해 프레임 간 급격한 스파이크를 제거합니다.
\end{enumerate}

출력은 엄격히 \textbf{운전자 권고(Advisory)} 용도로만 제공됩니다:

\begin{itemize}[leftmargin=1.5em,itemsep=2pt]
\item \texttt{Normal Driving} (정상 주행)
\item \texttt{Caution} (주의)
\item \texttt{Speed Reduction} (감속 권고)
\item \texttt{Braking Recommendation} (제동 권고)
\item \texttt{Emergency Intervention} (긴급 개입 권고)
\end{itemize}
\end{document}
